\documentclass[11pt,a4paper]{article}
\usepackage[T2A]{fontenc}
\usepackage[utf8]{inputenc}
\usepackage[russian]{babel}
\usepackage{amsmath,amssymb,amsthm,mathtools}
\usepackage[margin=2.3cm]{geometry}
\usepackage{enumitem}
\usepackage{tikz}
\usepackage[colorlinks=true,linkcolor=blue!50!black,citecolor=blue!50!black,urlcolor=blue!50!black]{hyperref}

\newtheorem{theorem}{Теорема}[section]
\newtheorem{lemma}[theorem]{Лемма}
\newtheorem{proposition}[theorem]{Предложение}
\newtheorem{corollary}[theorem]{Следствие}
\theoremstyle{definition}
\newtheorem{definition}[theorem]{Определение}
\newtheorem{example}[theorem]{Пример}
\newtheorem{remark}[theorem]{Замечание}

\newcommand{\R}{\mathbb R}

\newcommand{\eps}{\varepsilon}
\newcommand{\sym}{\operatorname{sym}}
\newcommand{\mk}[1]{\textsf{[#1]}}
\newcommand{\term}[1]{\emph{#1}}

\title{Неустойчивость равновесия в аналитическом потенциальном поле на плоскости\\[4pt]
\large Обзор работы В.\,П.~Паламодова (1977) для читателя вне теории устойчивости}
\author{Главный агент исследования (Claude) по поручению владельца}
\date{Черновик от 21.10.2026}

\begin{document}
\maketitle

\begin{abstract}
Пусть $U$ --- вещественно-аналитическая функция на плоскости, $0$ --- её критическая точка и не
точка минимума. Тогда положение равновесия $x=0$ системы $\ddot x=-\nabla U(x)$ неустойчиво.
Это доказал В.\,П.~Паламодов в 1977 году (независимо --- С.~Талиаферро, 1980). Текст объясняет,
что это значит и как доказывается: какие понятия механики здесь участвуют, откуда берётся
неустойчивость, где работает аналитичность и почему тот же путь пока не пройден в размерности
три и выше. Все специальные термины определены; их список с короткими пояснениями --- в
приложении~\ref{app:glossary}.
\end{abstract}

\paragraph{Состояние черновика.} Разделы \ref{sec:mech}--\ref{sec:method} и \ref{sec:after}
написаны. Раздел~\ref{sec:plane} (собственно конструкция Паламодова на плоскости) будет написан по
тексту статьи [P77]: с сервера, где работает исследование, сайт Math-Net.Ru недоступен, и текст
пока не прочитан. Там, где формулировки статьи известны лишь в пересказе других авторов, это
сказано явно.

\paragraph{Метки.} Каждое утверждение помечено. \mk{Литература: ссылка} --- опубликованный
результат; если первоисточник не открывался, это оговорено. \mk{Доказано} --- полное
доказательство приведено здесь. \mk{Проверено: кто, когда} --- доказательство независимо
прочитано рецензентом. Доказательства в этом тексте ещё ждут независимой проверки.

\tableofcontents

%=====================================================================
\section{Механика: движение, энергия, устойчивость}\label{sec:mech}

\subsection{Уравнение движения и фазовое пространство}

Точка (частица) единичной массы движется в $\R^n$. Её \term{положение} $x\in\R^n$, \term{скорость}
$\dot x=dx/dt$, ускорение $\ddot x$. На неё действует сила, заданная \term{потенциалом} (иначе:
\term{потенциальной энергией}) --- функцией $U\colon W\to\R$ на открытом множестве $W\subset\R^n$:
сила равна $-\nabla U(x)$. Закон Ньютона даёт
\begin{equation}\label{eq:newton}
\ddot x=-\nabla U(x).
\end{equation}
Пространство положений $\R^n$ (точнее, $W$) называют \term{конфигурационным пространством},
а $n$ --- \term{числом степеней свободы}. Уравнение второго порядка \eqref{eq:newton}
равносильно системе первого порядка
\begin{equation}\label{eq:system}
\dot x=v,\qquad \dot v=-\nabla U(x)
\end{equation}
на \term{фазовом пространстве} $W\times\R^n$, точки которого --- пары (положение, скорость).
\term{Движение} --- решение $x(t)$ уравнения \eqref{eq:newton}; \term{фазовая траектория} ---
кривая $t\mapsto(x(t),\dot x(t))$ в фазовом пространстве. Если $U\in C^2$, то правая часть
\eqref{eq:system} локально липшицева, и через каждую точку фазового пространства проходит
ровно одно решение (теорема Пикара); оно продолжается, пока не покидает компактные подмножества
$W\times\R^n$.

Такие системы называют \term{натуральными механическими системами}: энергия складывается из
кинетической $\tfrac12|\dot x|^2$ и потенциальной $U(x)$. В общей постановке кинетическая
энергия --- положительно определённая квадратичная форма скорости, зависящая от положения
($\tfrac12 g_x(\dot x,\dot x)$, где $g$ --- риманова метрика); у нас она евклидова. Предупреждение
о знаках: в части литературы используется \term{силовая функция} $-U$; тогда <<минимум
потенциала>> превращается в <<максимум силовой функции>>.

\subsection{Закон сохранения энергии}

\begin{lemma}\label{lem:energy}
\mk{Доказано} Функция $H(x,v)=\tfrac12|v|^2+U(x)$ (\term{полная энергия}) постоянна вдоль
каждого движения.
\end{lemma}
\begin{proof}
$\frac{d}{dt}H(x,\dot x)=\langle\dot x,\ddot x\rangle+\langle\nabla U(x),\dot x\rangle
=\langle\dot x,-\nabla U(x)+\nabla U(x)\rangle=0$.
\end{proof}

Значение $E=H$ на движении называют его \term{энергией}. Из $\tfrac12|\dot x|^2=E-U(x)\ge0$
следует, что движение с энергией $E$ не выходит из множества $\{x:U(x)\le E\}$ (\term{область
возможного движения}), а скорость в точке $x$ определяется положением: $|\dot x|=\sqrt{2(E-U(x))}$.
Это главный инструмент всего текста: если $E<0$, то движение целиком проходит там, где $U<0$.

\subsection{Положение равновесия и устойчивость по Ляпунову}

\term{Положение равновесия} --- постоянное движение $x(t)\equiv x_0$. По \eqref{eq:newton} это
возможно ровно тогда, когда $\nabla U(x_0)=0$, т.\,е. $x_0$ --- \term{критическая точка}
потенциала. Без ограничения общности $x_0=0$ и $U(0)=0$. Критическая точка \term{изолирована},
если в некоторой проколотой окрестности нуля $\nabla U\ne0$.

\begin{definition}[устойчивость по Ляпунову]\label{def:stab}
Равновесие $(x,v)=(0,0)$ \term{устойчиво по Ляпунову}, если для любого $\eps>0$ найдётся
$\delta>0$ такое, что каждое решение с $|x(0)|+|\dot x(0)|<\delta$ определено при всех $t\ge0$
и удовлетворяет $|x(t)|+|\dot x(t)|<\eps$ при всех $t\ge0$. Иначе равновесие
\term{неустойчиво}: существует $\eps>0$ такое, что при любом $\delta>0$ найдётся решение с
$|x(0)|+|\dot x(0)|<\delta$, которое в какой-то момент $t>0$ достигает $|x(t)|+|\dot x(t)|\ge\eps$.
\end{definition}

Иначе говоря, устойчивость означает: сколь угодно малое возмущение равновесия (по положению и
по скорости) остаётся малым навсегда. Неустойчивость: есть фиксированный масштаб $\eps$, до
которого раскачиваются некоторые сколь угодно малые возмущения.

\begin{example}
$n=1$. При $U=x^2/2$ (пружина) решения $x=a\cos t+b\sin t$ --- равновесие устойчиво.
При $U=-x^2/2$ (перевёрнутый маятник) решение $x=a\cosh t$ уходит от нуля --- неустойчиво.
\end{example}

Заметим: замена $t\mapsto-t$, $v\mapsto-v$ переводит решения в решения (\term{обратимость}),
поэтому всё сказанное о будущем верно и о прошлом.

\subsection{Теорема Лагранжа--Дирихле}

\begin{theorem}[Лагранж 1788, Дирихле 1846]\label{thm:LD}
\mk{Литература: классический результат; доказательство стандартное и приведено полностью,
\mk{Доказано}} Если $0$ --- точка \term{строгого локального минимума} $U$ (т.\,е. $U(x)>U(0)=0$
при $0<|x|\le r$), то равновесие устойчиво по Ляпунову.
\end{theorem}
\begin{proof}
Пусть $0<\eps\le r$ и $m(\eps)=\min_{|x|=\eps}U>0$. По непрерывности $H$ найдётся $\delta>0$
такое, что из $|x|+|v|<\delta$ следуют $|x|<\eps$ и $H(x,v)<m(\eps)$. Движение с такими
начальными данными не может достичь сферы $|x|=\eps$: там было бы $H\ge U\ge m(\eps)$. Значит,
$|x(t)|<\eps$ при всех $t$ (решение остаётся в компакте и продолжается на все $t$), а так как
$U\ge0$ в шаре, $|\dot x(t)|^2=2(H-U)<2m(\eps)$. Итак, $|x|+|\dot x|<\eps+\sqrt{2m(\eps)}$, и
это стремится к нулю при $\eps\to0$.
\end{proof}

Механический смысл: в ямке потенциала частица с малой энергией заперта стенками ямки.

\subsection{Обратная задача и задача Арнольда}

Естественный вопрос, восходящий к Ляпунову: верно ли обратное --- если $0$ не точка минимума $U$,
то равновесие неустойчиво? Нужно различать два случая отсутствия строгого минимума.

\begin{itemize}[leftmargin=*]
\item \term{Нестрогий минимум}: $U\ge0$ около нуля, но нули $U$ подходят к $0$. Пример:
$U=x_1^2$ на $\R^2$; тогда $x_1\equiv0$, $x_2=ct$ --- решение, и равновесие неустойчиво.
\item \term{Не минимум}: в любой окрестности нуля есть точки, где $U<0$.
\end{itemize}

Если критическая точка $0$ изолирована, нестрогий минимум невозможен: нуль $p\ne0$ функции
$U\ge0$ сам был бы точкой минимума, а значит критической точкой. \mk{Доказано}

\paragraph{Задача Арнольда 1971-4, случай (A).} В сборнике <<Задачи Арнольда>> задача 1971-4
цитируется (по [BMP21]; текст книги не сверен) так: <<Prove the instability of the equilibrium 0
of the analytic system $\ddot x=-\partial U/\partial x$ in the case where the isolated critical
point 0 of the potential U is not a minimum>>. То есть:

\medskip
\noindent\fbox{\parbox{0.96\linewidth}{\textbf{(A)} $U$ вещественно-аналитична около
$0\in\R^n$, $0$ --- изолированная критическая точка, не минимум. Доказать, что равновесие
неустойчиво по Ляпунову.}}
\medskip

Для $n=1$ это просто (см. ниже), для $n=2$ доказано Паламодовым [P77] и Талиаферро [Ta80];
при $n\ge3$ задача открыта \mk{Литература: Burgos--Paternain, arXiv:2208.01139, введение;
Burgos, J.~Differential Equations 454 (2026), arXiv:2506.19135, введение}.

\subsection{Почему важна аналитичность}\label{subsec:smooth}

\begin{example}[Пенлеве, Уинтнер]\label{ex:PW}
$n=1$, $U(x)=e^{-1/x^2}\cos(1/x)$, $U(0)=0$. \mk{Литература: Painlevé 1904; Wintner 1941; по
вторичным источникам. Проверка ниже --- \mk{Доказано}} Эта функция бесконечно
дифференцируема, но не аналитична в нуле (все её производные в нуле равны нулю, а сама она не
тождественный нуль). Точка $0$ не минимум: в точках $x=1/((2k+1)\pi)$ значение отрицательно. Но
равновесие устойчиво. Действительно, в точках $\pm a_k$, $a_k=1/(2k\pi)$, значения
$U(\pm a_k)=e^{-4k^2\pi^2}>0$ --- это <<барьеры>>. Частица на прямой с энергией
$E<U(\pm a_k)$, стартовавшая в $(-a_k,a_k)$, не может пройти через $\pm a_k$ (там было бы
$\tfrac12\dot x^2=E-U<0$). Если начальные данные малы, то и $E$ мало, и частица заперта в
сколь угодно малом интервале $(-a_k,a_k)$ со скоростью $|\dot x|\le\sqrt{2(E+\max|U|)}$. Здесь
критическая точка $0$ \emph{не} изолирована: между соседними нулями $U$ лежит критическая
точка (теорема Ролля), и они накапливаются к нулю.
\end{example}

\begin{proposition}\label{prop:n1}
\mk{Доказано} Случай (A) при $n=1$ верен даже для $U\in C^1$.
\end{proposition}
\begin{proof}
Так как $U'\ne0$ на $(0,r)$ и на $(-r,0)$, $U$ строго монотонна на каждом из этих интервалов.
Раз $0$ не минимум, на одном из них, скажем на $(0,r)$, $U$ убывает и отрицательна. Частица,
отпущенная с нулевой скоростью из $x_0\in(0,r)$, имеет $\ddot x=-U'(x)>0$, поэтому движется
вправо, и $\tfrac12\dot x^2=U(x_0)-U(x)$ растёт. Остановиться или стремиться к точке внутри
$(0,r)$ она не может: предельная точка была бы положением равновесия. Значит, она достигает
$x=r/2$ (при $x_0<r/2$), сколь бы мало ни было $x_0$.
\end{proof}

При $n\ge2$ гладкие примеры тоньше. Лалуа (1976) построил гладкий $U$ на плоскости, не имеющий
минимума в нуле, с устойчивым равновесием; там критическая точка тоже не изолирована.
Уренья (2020) построил гладкий $U$ на плоскости со строгим максимумом в нуле, $\nabla U\ne0$ вне
нуля, и с периодическими движениями сколь угодно близко к равновесию. Равновесие неустойчиво
(максимум), но не <<тотально>> (определение~\ref{def:total}). \mk{Литература: Laloy, Int.\ J.\
Non-Linear Mech.\ 11 (1976); Ureña, Ann.\ Mat.\ Pura Appl.\ 199 (2020); по вторичным
источникам} Вывод: в доказательстве (A) аналитичность (или изолированность) обязана где-то
работать по существу. Есть ли гладкий пример с \emph{изолированной} критической точкой без
минимума и с устойчивым равновесием --- неизвестно.

\subsection{Три способа доказать неустойчивость}\label{subsec:three}

\paragraph{1. Асимптотическая траектория.} Так называют движение $x(t)\not\equiv0$, для
которого $(x(t),\dot x(t))\to(0,0)$ при $t\to-\infty$ (или при $t\to+\infty$).

\begin{lemma}\label{lem:asympt}
\mk{Доказано} Если есть асимптотическая траектория, равновесие неустойчиво.
\end{lemma}
\begin{proof}
По обратимости можно считать, что $(x,\dot x)\to0$ при $t\to-\infty$. Положим
$\eps=|x(0)|+|\dot x(0)|>0$. Для любого $\delta>0$ найдётся $t_0<0$ с
$|x(t_0)|+|\dot x(t_0)|<\delta$; решение, стартующее в момент $t_0$, к моменту $0$ достигает
$\eps$.
\end{proof}

Этим путём шли Козлов, Талиаферро и другие: строили такое решение рядами по степеням $t$.

\paragraph{2. Функция Четаева.} Это функция $V(x,v)$ на фазовом пространстве, положительная в
некоторой области $D$, к границе которой примыкает равновесие, и возрастающая вдоль движений
в $D$ со скоростью, отделённой от нуля, пока движение близко к равновесию. Тогда движения из $D$,
стартующие сколь угодно близко к равновесию, уходят на фиксированное расстояние: иначе $V$
росла бы неограниченно, оставаясь ограниченной. Метод Паламодова --- частный случай этой идеи
с очень специальной функцией $V=\langle\dot x,v(x)\rangle$ (раздел~\ref{sec:method}).

\paragraph{3. Тотальная неустойчивость.}

\begin{definition}[тотальная неустойчивость]\label{def:total}
Критическая точка $0$ \term{тотально неустойчива}, если существуют окрестность $W_0$ нуля в
конфигурационном пространстве и функция $T\colon(-\infty,0)\to(0,\infty)$ такие, что каждое
движение с энергией $E<0$ проводит в $W_0$ без перерыва не больше времени $T(E)$: если
$x(t)\in W_0$ при $t\in[t_1,t_2]$, то $t_2-t_1\le T(E)$.
\end{definition}

\begin{remark}
\mk{Литература} В такой форме (<<с функцией $T(E)$>>) утверждение анонсировал Паламодов в 1995
году для всех $n$; так его цитирует Burgos (arXiv:2108.05829v4, введение). В обзоре Паламодова
2020 года [P20] определение формулируется иначе --- как уход из окрестности за конечное время
(по пересказу участника доски; первоисточник ещё не сверен). У Burgos (2025) определение другое
и слабее. Для нас важна форма определения~\ref{def:total}: её даёт метод.
\end{remark}

\begin{lemma}\label{lem:total-lyap}
\mk{Доказано} Если $0$ тотально неустойчива и не минимум, равновесие неустойчиво по Ляпунову.
Изолированность для этого не нужна.
\end{lemma}
\begin{proof}
Пусть шар $\{|x|<\eps\}$ лежит в $W_0$. Возьмём $x_0$ с $|x_0|<\delta$ и $U(x_0)<0$ (есть, так
как $0$ не минимум) и нулевую начальную скорость. Энергия движения $E=U(x_0)<0$, поэтому за
время $T(E)$ оно покидает $W_0$, а значит, достигает $|x|\ge\eps$.
\end{proof}

Тотальная неустойчивость сильнее, чем нужно: она говорит обо \emph{всех} движениях с
отрицательной энергией, а не об одном. Зато она хорошо подходит для доказательства <<одной
функцией>>, и именно её дают метод Паламодова и все результаты раздела~\ref{sec:method}.
\emph{Гипотеза Паламодова}: всякая критическая точка аналитического потенциала тотально
неустойчива \mk{Литература: [P20], гипотеза 1, по пересказу; не сверено}.

%=====================================================================
\section{Что нужно знать об аналитических функциях}\label{sec:analytic}

\subsection{Ряд Тейлора и младшая форма}

Функция $U$ \term{вещественно-аналитична} около $0$, если в некотором шаре она равна сумме
сходящегося степенного ряда. Группируя члены по степеням,
\[
U=u_m+u_{m+1}+u_{m+2}+\dots,
\]
где $u_k$ --- \term{однородная форма} степени $k$ (однородный многочлен: $u_k(\lambda x)=
\lambda^ku_k(x)$), а $u_m\not\equiv0$ --- \term{младшая форма}; $m$ --- \term{порядок} $U$ в нуле.
Если $0$ --- критическая точка и $U(0)=0$, то $m\ge2$. Форма $u_2$ задаётся \term{гессианом}
(матрицей вторых производных) $D^2U(0)$: $u_2(x)=\tfrac12\langle D^2U(0)x,x\rangle$.

\begin{itemize}[leftmargin=*]
\item Если $u_m$ принимает отрицательные значения, то и $U$ принимает: $U(tx)=t^m(u_m(x)+O(t))$.
\item Если $u_m>0$ вне нуля, то $0$ --- строгий минимум.
\item Трудный случай: $u_m\ge0$, но обращается в нуль на некоторых прямых. Тогда младшая форма
не решает дела. Пример: $U=y^2-x^4$ ($u_2=y^2$) --- не минимум; $U=y^2+x^4$ --- минимум.
\end{itemize}

Почти все классические результаты о неустойчивости (Ляпунов, Четаев, Хагедорн, Козлов,
Козлов--Паламодов) относятся к случаю, когда отсутствие минимума видно по младшей форме или
по другой <<главной части>> ряда. Ценность результата Паламодова для $n=2$ в том, что
на плоскости он не требует \emph{никаких} условий такого рода.

\subsection{Критические значения около нуля}

\begin{lemma}\label{lem:critval}
\mk{Доказано, по модулю леммы о выборе кривой} Пусть $U$ вещественно-аналитична около $0\in\R^n$
и $U(0)=0$. Тогда найдётся $\rho>0$ такое, что во всех критических точках $x\in B_\rho$
выполнено $U(x)=0$. Следовательно, для любого $E<0$
\[
\delta(E)=\min\{|\nabla U(x)|:\ |x|\le\rho/2,\ U(x)\le E\}>0
\]
(минимум по пустому множеству считаем равным $+\infty$).
\end{lemma}
\begin{proof}
Множество $C=\{\nabla U=0\}$ задаётся аналитическими уравнениями. Если бы критические точки
$x_k\to0$ имели $U(x_k)\ne0$, то $0$ лежал бы в замыкании множества $C\cap\{U\ne0\}$. По
\term{лемме о выборе кривой} \mk{Литература: Milnor, Singular points of complex hypersurfaces,
1968, лемма 3.1 (полуалгебраический случай); Łojasiewicz (полуаналитический случай); не сверено
с первоисточником} существует непрерывная кривая $\gamma\colon[0,1)\to\R^n$, аналитическая на
$(0,1)$, с $\gamma(0)=0$ и $\gamma(s)\in C\cap\{U\ne0\}$ при $s>0$. Но вдоль неё
$\frac{d}{ds}U(\gamma(s))=\langle\nabla U(\gamma(s)),\gamma'(s)\rangle=0$, так что
$U(\gamma(s))=U(0)=0$ --- противоречие. Второе утверждение: множество
$\{|x|\le\rho/2,\ U\le E\}$ компактно и не содержит критических точек ($U=E<0$ там, а критические
значения равны нулю), а $|\nabla U|$ непрерывна.
\end{proof}

Это первое место, где работает аналитичность. В примере Пенлеве--Уинтнера критические точки с
отрицательными значениями $U$ накапливаются к нулю, и лемма неверна.

\subsection{Нули и знаки аналитической функции на плоскости}\label{subsec:puiseux}

На плоскости устройство множества $\{U<0\}$ около нуля описывается полностью.

\begin{theorem}[Ньютон--Пюизё; локальная структура вещественной аналитической кривой]\label{thm:puiseux}
\mk{Литература: стандартный факт; см., например, C.~T.~C.~Wall, Singular points of plane curves,
2004, гл.~2; Brieskorn--Knörrer, Plane algebraic curves, §8.3; не сверено с первоисточником
в той форме, в какой оно здесь выписано} Пусть $U$ вещественно-аналитична около $0\in\R^2$,
$U(0)=0$, $U\not\equiv0$. Тогда при малом $\rho>0$ множество $\{U=0\}\cap B_\rho\setminus\{0\}$
--- объединение конечного числа $k\ge0$ попарно непересекающихся \term{полуветвей}
$\Gamma_1,\dots,\Gamma_k$. Каждая полуветвь пересекает каждую окружность $|x|=r$, $0<r<\rho$,
ровно в одной точке и в подходящих (повёрнутых) координатах задаётся как
\[
\Gamma_j=\{(x,y):\ 0<x<\rho_j,\ y=\textstyle\sum_{i\ge1}a_{ji}x^{i/N}\},
\]
где ряд по дробным степеням $x^{1/N}$ (\term{ряд Пюизё}) сходится, а повороты подобраны так, что
полуветвь касается в нуле положительной полуоси $x$ (в ряде нет членов степени $\le1$).
\end{theorem}

Отсюда: проколотый круг $B_\rho\setminus\{0\}$ разрезан полуветвями на $k$ \term{криволинейных
секторов} (при $k=0$ --- целиком один), и на каждом секторе $U$ сохраняет знак. Секторы, где
$U<0$, мы будем называть \term{рогами}. Если $0$ не минимум, есть хотя бы один рог. Рог может
иметь ненулевой угол раствора в нуле (например, квадрант у седла $U=xy$) или нулевой
(<<клюв>>), когда ограничивающие его полуветви касаются одной прямой; ширина рога на
расстоянии $r$ от нуля тогда порядка $r^\alpha$ с $\alpha>1$.

\begin{example}\label{ex:horns}
\begin{enumerate}[leftmargin=*]
\item $U=-(x^2+y^2)$ (строгий максимум): $k=0$, весь проколотый круг --- один <<рог>>.
\item $U=xy$ (седло): $k=4$ (полуоси), два рога --- второй и четвёртый квадранты.
\item $U=y^2-x^4$: полуветви $y=\pm x^2$ при $x>0$ и при $x<0$; два рога $\{|y|<x^2\}$
вдоль оси $x$, ширина $\sim r^2$.
\item $U=(y-x^2)^2-x^6$: рога $\{|y-x^2|<|x|^3\}$ --- узкие ($\sim r^3$) и \emph{искривлённые}:
они идут вдоль параболы $y=x^2$. Нуль --- изолированная критическая точка
($\partial_yU=0\Rightarrow y=x^2\Rightarrow\partial_xU=-6x^5=0$).
\end{enumerate}
\end{example}

\begin{figure}[ht]
\centering
\begin{tikzpicture}[scale=2.6]
\draw[->,gray] (-0.1,0)--(1.5,0) node[right]{$x$};
\draw[->,gray] (0,-0.1)--(0,1.9) node[above]{$y$};
\fill[blue!15] plot[domain=0:1.35,samples=60,smooth] (\x,{\x*\x+0.2*\x*\x*\x}) -- plot[domain=1.35:0,samples=60,smooth] (\x,{\x*\x-0.2*\x*\x*\x}) -- cycle;
\draw[blue] plot[domain=0:1.35,samples=60,smooth] (\x,{\x*\x+0.2*\x*\x*\x});
\draw[blue] plot[domain=0:1.35,samples=60,smooth] (\x,{\x*\x-0.2*\x*\x*\x});
\draw[dashed,gray] plot[domain=0:1.32,samples=40,smooth] (\x,{\x*\x});
\fill[red] (1,0.9) circle (0.02);
\draw[->,red,thick] (1,0.9) -- (1.33,1.197);
\draw[dotted,red] (0,0) -- (1,0.9);
\node[red,right] at (1.33,1.15) {\small радиальное направление};
\node at (0.35,1.0) {\small $U>0$};
\node at (1.15,0.45) {\small $U>0$};
\draw[blue!60!black] (0.62,0.39) -- (0.95,0.2) node[right]{\small рог $U<0$};
\end{tikzpicture}
\caption{Искривлённый рог вдоль параболы $y=x^2$ (пунктир), как у функции $U=(y-x^2)^2-x^6$;
ширина рога для наглядности преувеличена. Из точки нижней половины рога радиальное направление
(красная стрелка, продолжение луча из нуля) идёт положе параболы и выходит из рога через нижнюю
границу, туда, где $U>0$.}
\label{fig:horn}
\end{figure}

%=====================================================================
\section{Метод Паламодова: функция $\langle\dot x,v(x)\rangle$}\label{sec:method}

\subsection{Образец: однородный потенциал и вириал}

Пусть $U$ однороден степени $m\ge2$ ($U(\lambda x)=\lambda^mU(x)$) и где-то отрицателен.
Рассмотрим \term{вириал} $I(t)=\tfrac12|x(t)|^2$ --- половину квадрата расстояния до нуля.
По тождеству Эйлера $\langle x,\nabla U(x)\rangle=mU(x)$, поэтому
\[
\ddot I=|\dot x|^2+\langle x,\ddot x\rangle=|\dot x|^2-mU=2(E-U)-mU=2E-(m+2)U
\]
(\term{тождество Лагранжа--Якоби}). Если $E<0$, то $U\le E$ на движении, и
\[
\ddot I\ \ge\ 2E-(m+2)E=m|E|>0.
\]

\begin{proposition}\label{prop:homog}
\mk{Доказано} Для однородного $U$, принимающего отрицательные значения, равновесие тотально
неустойчиво и неустойчиво по Ляпунову.
\end{proposition}
\begin{proof}
Если движение с энергией $E<0$ проводит в шаре $|x|<\rho$ время $[0,T]$, то
$|\dot I|=|\langle x,\dot x\rangle|\le\rho\sqrt{2(E-\min_{B_\rho}U)}=:\rho w$, и из
$\ddot I\ge m|E|$ получаем $m|E|T\le\dot I(T)-\dot I(0)\le2\rho w$, т.\,е. $T\le2\rho w/(m|E|)$.
Это тотальная неустойчивость; так как $U(\lambda x_0)=\lambda^mU(x_0)<0$ при всех $\lambda>0$,
нуль не минимум, и работает лемма~\ref{lem:total-lyap}.
\end{proof}

Здесь $\dot I=\langle\dot x,x\rangle$ --- это функция Четаева вида $\langle\dot x,v(x)\rangle$ с
полем $v(x)=x$. Свойства поля, которые сработали:
\begin{enumerate}[label=(\roman*)]
\item поле <<толкает вниз по потенциалу>>: $\langle v,\nabla U\rangle=mU<0$ на $\{U<0\}$;
\item поле <<растягивает>>: производная $Dv=I$ положительно определена.
\end{enumerate}
Для неоднородного $U$ радиальное поле ломается (пример~\ref{ex:radial-fails}), и вся идея
Паламодова --- найти поле со свойствами (i), (ii) для любого аналитического $U$ на плоскости.

\subsection{Критерий поля}

Для квадратной матрицы $M$ через $\sym M=\tfrac12(M+M^{\mathsf T})$ обозначим её симметричную
часть; $\sym M\succeq c$ значит $\langle M\xi,\xi\rangle\ge c|\xi|^2$ при всех $\xi$. Поле $v$ с
$\sym Dv\succeq0$ называют \term{монотонным}, с $\sym Dv\succeq c>0$ --- \term{сильно
монотонным}: тогда $\langle v(x)-v(y),x-y\rangle\ge c|x-y|^2$ на выпуклых множествах.
Пример: градиент выпуклой функции.

Ключевое вычисление. Пусть $v$ --- $C^1$-поле и $x(t)$ --- движение. Для
$F(t)=\langle\dot x(t),v(x(t))\rangle$
\begin{equation}\label{eq:Fdot}
\dot F=\langle\ddot x,v\rangle+\langle\dot x,Dv\,\dot x\rangle
=-\langle\nabla U,v\rangle+\langle\sym Dv\,\dot x,\dot x\rangle .
\end{equation}
Первый член неотрицателен, если поле толкает вниз; второй --- если поле монотонно.

\begin{theorem}[критерий поля]\label{thm:field}
\mk{Доказано; то же рассуждение проверено рецензентом: reviewer R1, 21.10.2026,
\texttt{notes/reviewer/2026-10-21-max-gluing.md} (лемма~1 из
\texttt{notes/vlad-claude/theoremB-proof.md})} Пусть $U\in C^2$ около замкнутого шара
$\bar B_\rho$, $U(0)=0$, $X=\{x\in B_\rho:U(x)<0\}$. Пусть на открытом множестве, содержащем
$X$, задано $C^1$-поле $A$ с $|A|\le M$ на $X$ и
\[
\sym DA(x)\succeq0,\qquad \langle A(x),\nabla U(x)\rangle\le p\,U(x)\qquad(x\in X)
\]
при некотором $p>0$. Тогда любое движение с энергией $E<0$, остающееся в $B_\rho$ при
$t\in[0,T]$, удовлетворяет $T\le 2wM/(p|E|)$, где $w=\sqrt{2(E-\min_{\bar B_\rho}U)}$. В
частности, $0$ тотально неустойчива.
\end{theorem}
\begin{proof}
На движении $U(x(t))\le E<0$, т.\,е. $x(t)\in X$, и $|\dot x|\le w$. По \eqref{eq:Fdot}
$\dot F\ge-pU(x)\ge p|E|$, а $|F|\le wM$. Значит, $2wM\ge F(T)-F(0)\ge p|E|T$.
\end{proof}

Паламодов в [P77] использует вариант, где условия переставлены: толчок слабее
($\langle v,\nabla U\rangle\le0$), растяжение сильнее ($\sym Dv\succeq c>0$), а поле может быть
лишь непрерывным и гладким вне конечного числа кривых. Точная формулировка его основной леммы
будет выписана по тексту [P77]; по пересказу участника доски, прочитавшего начало статьи, она
такова: \emph{если на $U^-_\rho=\{|x|<\rho,\ U<0\}$ есть непрерывное поле $v$, гладкое вне
конечного числа кривых, с $\langle v,\nabla U\rangle\le0$ и $\langle Dv\,\xi,\xi\rangle\ge c|\xi|^2$,
то равновесие неустойчиво} \mk{Литература: [P77], основная лемма; не сверено с первоисточником}.
Вот почему такой вариант сводится к теореме~\ref{thm:field}, когда $U$ аналитична.

\begin{proposition}\label{prop:palamodov-variant}
\mk{Доказано} Пусть $U$ вещественно-аналитична около $0$, $U(0)=0$, и $\rho$ взято из
леммы~\ref{lem:critval}. Пусть на открытом множестве, содержащем $X=\{|x|<\rho/2,\ U<0\}$,
задано ограниченное $C^1$-поле $v$ с $\langle v,\nabla U\rangle\le0$ и $\sym Dv\succeq c>0$ на
$X$. Тогда любое движение с энергией $E<0$, остающееся в шаре $|x|<\rho/2$ при $t\in[0,T]$,
удовлетворяет $T\le 2w(M+\sigma L)/(\sigma\,\delta(E)^2)$, где $M=\sup_X|v|$,
$L=\sup_X|\nabla U|$, $\sigma=c/\sup_{X}\lVert D^2U\rVert$, а $w$ и $\delta(E)$ --- как выше.
В частности, $0$ тотально неустойчива.
\end{proposition}
\begin{proof}
Положим $v_\sigma=v-\sigma\nabla U$. Тогда $\sym Dv_\sigma=\sym Dv-\sigma D^2U\succeq
c-\sigma\lVert D^2U\rVert\succeq0$ и $\langle v_\sigma,\nabla U\rangle\le-\sigma|\nabla U|^2$. На
движении с энергией $E<0$ точка $x(t)$ лежит в $\{|x|<\rho/2,\ U\le E\}$, где
$|\nabla U|\ge\delta(E)>0$ (лемма~\ref{lem:critval}). По \eqref{eq:Fdot} для
$F=\langle\dot x,v_\sigma(x)\rangle$ имеем $\dot F\ge\sigma\delta(E)^2$, а $|F|\le w(M+\sigma L)$.
\end{proof}

\begin{remark}\label{rem:where-analytic}
Где здесь аналитичность: только в лемме~\ref{lem:critval}, т.\,е. в том, что на подуровне
$\{U\le E\}$ около нуля нет критических точек. Вся остальная трудность --- \emph{построить}
поле. Для гладкого потенциала Уреньи такого поля нет (иначе равновесие было бы тотально
неустойчиво), так что именно построение поля обязано использовать аналитичность ещё раз.
\end{remark}

\begin{remark}[непрерывные поля]\label{rem:lipschitz}
\mk{Доказано, эскиз проверки} Гладкость поля нужна лишь для формулы \eqref{eq:Fdot}. Если $v$
непрерывно и кусочно гладко (гладко вне конечного числа кривых) и на каждом куске
$\sym Dv\succeq c$, то $v$ сильно монотонно на каждом выпуклом подмножестве (интегрированием
вдоль отрезка, пересекающего кривые в конечном числе точек), и для липшицевой функции
$F(t)=\langle\dot x,v(x(t))\rangle$ почти всюду
$\dot F\ge-\langle\nabla U,v\rangle+c|\dot x|^2$: второй член --- предел
$\langle\frac{x(t+h)-x(t)}h,\frac{v(x(t+h))-v(x(t))}h\rangle\ge c\,\bigl|\frac{x(t+h)-x(t)}h\bigr|^2$.
Поэтому кусочно гладкие непрерывные поля годятся, и склеивать поля можно <<без сглаживания>>,
лишь бы склеенное поле оставалось непрерывным. Это существенная свобода: при сглаживании
разбиением единицы возникает член $(v_1-v_2)\otimes\nabla\chi$, который портит монотонность
(именно на нём споткнулось отозванное доказательство Burgos 2021 для произвольного $n$).
\end{remark}

\subsection{Почему радиальное поле не годится}

\begin{example}\label{ex:radial-fails}
\mk{Доказано, прямое вычисление} Для $U=(y-x^2)^2-x^6$ (пример~\ref{ex:horns}.4) положим
$s=y-x^2$. Для радиального поля $v=(x,y)$:
\[
\langle v,\nabla U\rangle=x(-4xs-6x^5)+2ys=2s^2-2x^2s-6x^6 .
\]
В нижней половине рога, $s=-\lambda x^3$ с $0<\lambda<1$, это равно $2\lambda x^5+(2\lambda^2-6)x^6>0$
при малых $x>0$: радиальное поле толкает <<вверх по потенциалу>>, поперёк искривлённого рога
(рис.~\ref{fig:horn}). Зато поле $v=(x,2y)$, касающееся парабол $y=cx^2$, годится:
\[
\langle v,\nabla U\rangle=-4x^2s-6x^6+4ys=4s^2-6x^6=4U-2x^6\le4U,\qquad Dv=\operatorname{diag}(1,2),
\]
и по теореме~\ref{thm:field} равновесие тотально неустойчиво.
\end{example}

Мораль: поле надо подстраивать под геометрию каждого рога. Поле вида
$v=(x,\,w\,y)$ --- \term{взвешенное поле Эйлера}; его поток $x\mapsto e^tx$, $y\mapsto e^{wt}y$
сохраняет кривые $y=cx^w$. Если в рогу $U$ <<почти квазиоднороден>> относительно весов $(1,w)$
(т.\,е. $U(\lambda x,\lambda^wy)\approx\lambda^dU(x,y)$), то такое поле работает. Но в одном
рогу разные его части могут требовать разных весов: например, у $U=(y-x^2)(y-x^3)$ рог
$x^3<y<x^2$ у верхнего края похож на кривую веса $2$, у нижнего --- веса $3$, а посередине
$U\approx-x^2y$. Тогда ни одно поле $(x,wy)$ не годится \mk{Доказано, прямая проверка: при
каждом $w$ и $p$ неравенство $\langle v,\nabla U\rangle\le pU$ нарушается где-то в рогу}, и поля
разных весов приходится склеивать. Как Паламодов делает это на плоскости --- содержание
раздела~\ref{sec:plane}.

%=====================================================================
\section{Конструкция поля на плоскости (по [P77])}\label{sec:plane}

\emph{Раздел будет написан по тексту статьи [P77].} Что известно из вторичных источников:

\begin{itemize}[leftmargin=*]
\item Результат: для вещественно-аналитического $U$ на $\R^2$ критическая точка, не являющаяся
минимумом, --- неустойчивое равновесие \mk{Литература: Burgos--Paternain, arXiv:2208.01139,
введение; Enciso--Peralta-Salas, arXiv:2610.27893}. Изолированность не предполагается. По
[P20], в [P77] доказана даже тотальная неустойчивость \mk{по пересказу; не сверено}.
\item Метод: основная лемма (раздел~\ref{sec:method}) и построение поля $v$ на каждом роге.
По пересказу участника доски, двумерная конструкция использует угол между векторами $-x$ и
$\nabla U(x)$ \mk{не сверено}. По тому же пересказу, поле строится и для бесконечно
дифференцируемых $U$ с критической точкой конечной кратности \mk{не сверено}.
\end{itemize}

%=====================================================================
\section{Что было дальше и почему не получается при $n\ge3$}\label{sec:after}

\paragraph{Другие доказательства при $n=2$.}
\begin{itemize}[leftmargin=*]
\item Талиаферро (1980) независимо доказал неустойчивость по Ляпунову для аналитического
$U$ на плоскости \mk{Литература: S.\,D.~Taliaferro, J.~Differential Equations 35 (1980),
248--265; по вторичным источникам}.
\item Козлов (1981) перенёс результат на произвольную кинетическую энергию, опираясь на работу
Паламодова \mk{Литература: В.\,В.~Козлов, УМН 36:1 (1981), 209--210; по вторичным источникам}.
\item Лалуа и Пайфер (1982) разобрали нестрогие минимумы на плоскости; вместе с предыдущими это
даёт полное обращение теоремы Лагранжа--Дирихле для аналитических $U$ при $n=2$
\mk{Литература: Burgos--Paternain, arXiv:2208.01139, введение}.
\item Брюнелла (1998) дал другое, топологическое доказательство (A) при $n=2$: при $n=2$ уровень
энергии $\{H=0\}$ трёхмерен, и его теорема о векторных полях в размерности три даёт
асимптотическую траекторию \mk{Литература: M.~Brunella, Ann.\ Inst.\ Fourier 48 (1998),
1345--1357; формулировку сверял агент по литературе}.
\end{itemize}

\paragraph{Попытки для всех $n$.} В 1995 году Паламодов анонсировал тотальную неустойчивость
для всех $n$; полного доказательства опубликовано не было \mk{Литература: Allaire--Rauch,
ARMA 224 (2017), §1.7; Burgos, arXiv:2108.05829v4, введение}. В обзоре 2020 года это
<<гипотеза 1>> \mk{по пересказу; не сверено}. Препринт Burgos (2021) с доказательством для всех
$n$ отозван автором: ошибка в лемме~2.2 о склейке локальных полей \mk{Литература: arXiv:2108.05829,
комментарий к v5}. В работе Burgos (2025) тотальная неустойчивость доказана для <<общих>>
потенциалов (когда младшая форма на единичной сфере --- функция Морса с регулярным уровнем $0$),
а случай (A) назван открытым.

\paragraph{Что ломается при $n\ge3$.} На плоскости множество $\{U<0\}$ около нуля --- конечное
число рогов, каждый ограничен двумя кривыми с явными разложениями Пюизё. В $\R^3$ это области,
ограниченные поверхностями; их сечения сферой $|x|=r$ --- области на двумерной сфере с
нетривиальной топологией, а явной параметризации нет. Нужны раздутия (разрешение особенностей),
а при переходе через раздутие монотонность поля ($\sym Dv\succeq0$) не сохраняется, и склейка
локальных полей вносит член без нужной малости. Это и есть открытая часть задачи.

%=====================================================================
\appendix
\section{Словарь терминов}\label{app:glossary}

\begin{description}[leftmargin=0pt,style=nextline]
\item[Асимптотическая траектория] движение, стремящееся (вместе со скоростью) к равновесию при
$t\to+\infty$ или $t\to-\infty$; её существование влечёт неустойчивость (лемма~\ref{lem:asympt}).
\item[Вириал] $I=\tfrac12|x|^2$; его вторая производная по времени --- тождество Лагранжа--Якоби.
\item[Взвешенное поле Эйлера] $v=(w_1x_1,\dots,w_nx_n)$; для квазиоднородного $U$ веса $(w_i)$
и степени $d$ выполнено $\langle v,\nabla U\rangle=dU$.
\item[Квазиоднородная функция] $U(\lambda^{w_1}x_1,\dots,\lambda^{w_n}x_n)=\lambda^dU(x)$ при всех
$\lambda>0$.
\item[Критическая точка] точка, где $\nabla U=0$; изолированная --- если других рядом нет.
\item[Лемма о выборе кривой] если точка лежит в замыкании полуаналитического множества, к ней
ведёт аналитическая дуга из этого множества.
\item[Младшая форма] первая ненулевая однородная часть ряда Тейлора.
\item[Монотонное поле] поле $v$ с $\sym Dv\succeq0$; сильно монотонное --- $\succeq c>0$.
\item[Натуральная механическая система] система с энергией <<кинетическая (квадратичная по
скоростям) плюс потенциальная (зависит только от положения)>>.
\item[Неустойчивость по Ляпунову] отрицание устойчивости (определение~\ref{def:stab}).
\item[Полуветвь] одна из конечного числа дуг, на которые распадается $\{U=0\}$ около нуля на
плоскости (теорема~\ref{thm:puiseux}).
\item[Потенциал, потенциальная энергия] функция $U$; сила равна $-\nabla U$. Силовая функция ---
$-U$.
\item[Равновесие (положение равновесия)] постоянное решение; при $\dot x=v$, $\dot v=-\nabla U$
--- точка $(x_0,0)$ с $\nabla U(x_0)=0$.
\item[Рог] связная компонента множества $\{U<0\}$ в малой проколотой окрестности нуля.
\item[Ряд Пюизё] сходящийся ряд по дробным степеням $x^{1/N}$.
\item[Строгий, нестрогий минимум] строгий: $U>U(0)$ в проколотой окрестности; нестрогий:
$U\ge U(0)$, но не строгий.
\item[Тотальная неустойчивость] все движения с энергией ниже $U(0)$ покидают фиксированную
окрестность за время, ограниченное через энергию (определение~\ref{def:total}).
\item[Устойчивость по Ляпунову] малые возмущения положения и скорости остаются малыми навсегда
(определение~\ref{def:stab}).
\item[Фазовое пространство] пространство пар (положение, скорость).
\item[Функция Четаева] функция на фазовом пространстве, растущая вдоль движений вблизи
равновесия в области, примыкающей к нему; доказывает неустойчивость.
\item[$\sym M$] симметричная часть матрицы, $\tfrac12(M+M^{\mathsf T})$.
\end{description}

%=====================================================================
\begin{thebibliography}{99}
\bibitem[P77]{P77} В.\,П.~Паламодов, Об устойчивости равновесия в потенциальном поле,
\emph{Функц. анализ и его прил.} 11:4 (1977), 42--55; англ.: \emph{Funct. Anal. Appl.} 11
(1977), 277--289, doi:10.1007/BF01077142.
\bibitem[P95]{P95} V.\,P.~Palamodov, Stability of motion and algebraic geometry, in:
\emph{Dynamical Systems in Classical Mechanics} (V.\,V.~Kozlov, ed.), Amer. Math. Soc. Transl.
Ser.~2, vol.~168, 1995, 5--20.
\bibitem[P20]{P20} В.\,П.~Паламодов, Об обращении теоремы Лагранжа--Дирихле и неустойчивости
консервативных систем, \emph{УМН} 75:3 (2020), 107--122.
\bibitem[Ta80]{Ta80} S.\,D.~Taliaferro, Stability for two dimensional analytic potentials,
\emph{J. Differential Equations} 35 (1980), 248--265.
\bibitem[BMP21]{BMP21} J.\,M.~Burgos, E.~Maderna, M.~Paternain, On the Lyapunov instability in
Newtonian dynamics, \emph{Nonlinearity} 34 (2021), 6719--6726; arXiv:1906.11962.
\bibitem[BP22]{BP22} J.\,M.~Burgos, M.~Paternain, On the Lagrange--Dirichlet converse in
dimension three, arXiv:2208.01139 (2022), препринт.
\bibitem[B21]{B21} J.\,M.~Burgos, A proof of the Palamodov's total instability conjecture,
arXiv:2108.05829 (2021; отозвано автором в 2022).
\bibitem[B25]{B25} J.\,M.~Burgos, McGehee blowup for Lagrangian systems and instability of
equilibria, \emph{J. Differential Equations} 454 (2026); arXiv:2506.19135.
\bibitem[Br98]{Br98} M.~Brunella, Instability of equilibria in dimension three, \emph{Ann. Inst.
Fourier} 48 (1998), 1345--1357.
\bibitem[EPS26]{EPS26} A.~Enciso, D.~Peralta-Salas, Stability and instability for particle
dynamics in magnetic fields, arXiv:2610.27893 (2026).
\end{thebibliography}

\end{document}
